% Clase 7 --- El experimento de la caja (§4).
% Al tocar el fondo, bolita y caja pasan a ser UN solo conductor: la carga
% migra a la superficie EXTERNA (§2.1) y la bolita queda con carga cero. Medir
% "cero o no cero" es muchísimo más preciso que medir una fuerza.
\begin{center}
\begin{tikzpicture}[scale=1]

  % =============== (1) bolita cargada, colgada ===============
  \begin{scope}
    \draw[figgray, line width=1.2pt] (-1.3,1.5) -- (-1.3,-1.3) -- (1.3,-1.3) -- (1.3,1.5);
    \draw[figgray, line width=1.2pt] (-1.3,1.5) -- (-0.15,1.5);
    \draw[figgray, line width=1.2pt] (0.15,1.5) -- (1.3,1.5);
    \draw[figaux] (0,1.5) -- (0,0.25);
    \node[figlblaux, anchor=south] at (0.0,1.62) {hilo aislante};
    \node[figpos, minimum size=10pt] at (0,0) {};
    \node[figlbl, figred, anchor=east] at (-0.2,0) {$q$};
    \node[figlbl] at (0,-2.15) {(1) bolita cargada};
  \end{scope}

  % =============== (2) toca el fondo ===============
  \begin{scope}[xshift=4.6cm]
    \draw[figgray, line width=1.2pt] (-1.3,1.5) -- (-1.3,-1.3) -- (1.3,-1.3) -- (1.3,1.5);
    \draw[figgray, line width=1.2pt] (-1.3,1.5) -- (-0.15,1.5);
    \draw[figgray, line width=1.2pt] (0.15,1.5) -- (1.3,1.5);
    \draw[figaux] (0,1.5) -- (0,-1.05);
    \node[figneutra, minimum size=10pt] at (0,-1.15) {};
    % la carga migró a la cara externa
    \foreach \x in {-1.05,-0.6,...,1.1} {\node[figlbl,figred,scale=0.65] at (\x,-1.5){$+$};}
    \foreach \y in {-1.0,-0.5,...,1.4} {
      \node[figlbl,figred,scale=0.65] at (-1.5,\y){$+$};
      \node[figlbl,figred,scale=0.65] at (1.5,\y){$+$};
    }
    \node[figlbl, align=center] at (0,-2.15) {(2) toca el fondo:\\ un solo conductor};
  \end{scope}

  % =============== (3) se mide la bolita ===============
  \begin{scope}[xshift=9.6cm]
    \node[figneutra, minimum size=10pt] at (0,0) {};
    \draw[figaux] (0,1.5) -- (0,0.25);
    \node[figlbl, figred, anchor=west] at (0.35,0) {$q = 0$};
    \draw[figvecaux] (-1.5,0) -- (-0.35,0);
    \node[figlblaux, anchor=east] at (-1.55,0) {se saca};
    \node[figlbl, align=center] at (0,-2.15) {(3) se mide:\\ carga \textbf{cero}};
  \end{scope}

\end{tikzpicture}\\[0.5em]
{\small\itshape Si la ley fuese $1/r^{2+\delta}$ con $\delta \neq 0$, la carga
\textbf{no} migraría por completo a la cara externa y la bolita saldría con algo
de carga. Como el experimento sólo pregunta \textbf{¿hay carga o no?}, y eso se
detecta muchísimo mejor que una fuerza, se llega a $|\delta| < 10^{-15}$. Lo
propuso \textbf{Priestley en 1767}, antes que Coulomb.}
\end{center}
